\documentclass[aspectratio=169,9pt]{beamer}
\usepackage[french]{babel}
\selectlanguage{french}
\usepackage[utf8]{inputenc}
\usepackage{framed}
\usetheme{unisante}

\usepackage{tikz, tikz-cd}
\usepackage{pgf}
\usepackage{graphicx}
\usepackage[font=small,labelfont=bf]{caption}
\usepackage{subcaption}
\usepackage{hyperref}
\usepackage{multirow}
\usepackage{listings}
\usepackage{amsmath}
\usepackage{qrcode}

%%% Compact itemize, mirroring \compresslist from the poster %%%%%%%%%%%%%%%%%%%
\newcommand{\compresslist}{%
  \setlength{\itemsep}{2pt}%
  \setlength{\parskip}{0pt}%
  \setlength{\parsep}{0pt}%
}
\newcommand{\compressl}{%
  \setlength{\itemsep}{4pt}%
  \setlength{\parskip}{2pt}%
  \setlength{\parsep}{2pt}%
}
\definecolor{monvert}{rgb}{0.0, 0.5, 0.0}
\newcommand{\smallitem}[1]{{\fontsize{8pt}{9pt}\selectfont #1}}
\setbeamerfont{footnote mark}{size=\tiny}
\setbeamerfont{footnote}{size=\tiny}

\title[]{Caractérisation de la baisse des naissances en Suisse : analyse statistique de 2000 à 2025}
\author{Anthony Hauser, Julien Riou, Kaspar Staub}
\institute{\normalsize Unisanté}
\date{\normalsize Lausanne, 10 septembre 2026}

%% Contenu
\begin{document}
	
	
	{
		\setbeamertemplate{footline}{} 
		\begin{frame}
			\titlepage
		\end{frame}
	}
	\addtocounter{framenumber}{-1}
	
	\section*{Introduction}

%%% Background %%%%%%%%%%%%%%%%%%%%%%%%%%%%%%%%%%%%%%%%%%%%%%%%%%%%%%%%%%%%%%%%
\begin{frame}{La baisse de la fécondité : un phénomène global}
\begin{minipage}[t]{0.4\linewidth}
\textbf{Baisse de la natalité}
\begin{itemize}\compresslist
\item \textbf{Indice conjoncturel de fécondité} (ICF) : nombre moyen d'enfants par femme, mesuré sur une année.\footnote[frame]{Peeples L., \emph{Nature} 2025;644(8077):594--596.}
\item Déclin de la fécondité sur \textbf{tous les continents}
\item Dans de nombreux pays, l'ICF est désormais \textbf{inférieur au niveau de remplacement} ($\sim$2.1 enfants par femme)
\item \textbf{Europe} : ICF $\approx$ 1.4
\end{itemize}
\end{minipage}
\begin{minipage}[t]{0.56\linewidth}
\textbf{Raisons du déclin}
\begin{itemize}\compresslist
\item \textbf{Report de la parentalité}: les femmes ont des enfants plus tard (\emph{tempo effect}).
\item Mais aussi une \textbf{réduction de la fécondité au cours de la vie} (\emph{quantum effect}) : les femmes ont moins d'enfants.\footnote[frame]{ Hellstrand et al., \emph{Demography} 2021;58(4):1373--1399.}
%étude par cohorte de femmes en âge de procréer, permet d'ajuster sur l'effect tempo et d'isoler l'effet quantum
\end{itemize}
\vspace*{0.2cm}
\includegraphics[width=1\linewidth,keepaspectratio]{pictures/birth_decline_nature.png}
\end{minipage}
\end{frame}



\begin{frame}{La baisse de la fécondité en Suisse}
\begin{minipage}[t]{0.48\linewidth}
\textbf{Evolution 2000-2025}
\begin{itemize}\compresslist
\item La baisse actuelle s'inscrit dans une \textbf{tendance de long terme}, déjà visible avant la pandémie.
\item En Suisse, l'ICF atteint \textbf{1,29 en 2025}, son plus bas niveau historique.
\item Exception: année 2021 où les naissances ont augmenté.
%COVID-19 : un \textbf{« pandemic roller-coaster »}, avec une baisse initiale, un rebond en 2021, puis une nouvelle baisse à partir de 2022.

\end{itemize}
\vspace*{0.2cm}
\centering
\includegraphics[width=0.8\linewidth,keepaspectratio]{pictures/icf_switzerland_2020_2025.png}
\end{minipage}
\hfill
\begin{minipage}[t]{0.48\linewidth}
\textbf{Comparaison historique}
\begin{itemize}\compresslist
\item Historiquement, les pandémies (1889–90, 1918–20, 1957) coincidaient avec un \textbf{déclin des naissances}, généralement suivies d'un \textbf{rebond}.\footnote[frame]{Matthes et al., \emph{Population Studies} 2025;79(3):463–478.}
\item On observe aussi que la phase de déclin observé actuellement n'a rien de nouveau.
\end{itemize}%permet de remettre en perspective
\vspace*{0.2cm}
\centering
\includegraphics[width=1\linewidth,keepaspectratio]{pictures/matthes_historical_birth.png}
\end{minipage}
\end{frame}


\begin{frame}{Une baisse de la natalité hétérogène}

\begin{minipage}[h!]{0.48\linewidth}
\textbf{Excès/Déficit de la natalité en Suisse}
\begin{itemize}\compresslist
\item Définition: Natalité observée -- Natalité attendue
\item \textbf{Des différences selon les populations}\footnote[frame]{Matthes et al., \emph{Population Studies} 2025;79(3):463–478.}
\begin{itemize}\compresslist
\item Région
\item Nationalité
\item Parité
\item Âge maternel
\end{itemize}
\end{itemize}
\end{minipage}
\hfill
\begin{minipage}[h!]{0.48\linewidth}
\textbf{Des analyses encore limitées}
\begin{itemize}\compresslist
\item Analyse \textbf{par sous-groupe} : tendances décrites séparément, sans \textbf{comparaison directe} des déficits
\item Ne prend pas en compte \textbf{l'âge maternel} % ce qui peut expliquer certaines différences observées par sous-groupe
\begin{itemize}\compresslist 
\item Différences de \textbf{structure par âge}
\item \textbf{Report de la parentalité}: augmentation de l'âge à la maternité
\end{itemize}
\end{itemize}
\end{minipage}

\vfill

\centering
\includegraphics[width=0.9\linewidth,keepaspectratio]{pictures/matthes_excess_subgroup.png}
\end{frame}

\begin{frame}{Objectifs}

\begin{itemize}\compresslist
\item \textbf{Quantifier l’évolution de la natalité en Suisse (2000–2025)} en tenant compte de la structure de la population par âge et du report de la parentalité.

\item \textbf{Explorer les différences géographiques} de natalité et leurs déterminants.
\end{itemize}
\end{frame}

\begin{frame}{Données}

\begin{minipage}[t]{0.31\linewidth}
\textbf{Naissances}\
\textbf{(BEVNAT, OFS)}
\begin{itemize}\compresslist
\item Données individuelles, \textbf{2000--2025}
\item Naissances vivantes, mères \textbf{15--50 ans}
\item Âge, commune de résidence, nationalité, parité, mois et année
\end{itemize}
\end{minipage}
\hfill
\begin{minipage}[t]{0.31\linewidth}
\textbf{Population}\
\textbf{(OFS)}
\begin{itemize}\compresslist
\item Femmes 15--50 ans
\item Données par \textbf{âge, année, nationalité et commune}
\item 2000--2025 selon le niveau de stratification
\end{itemize}
\end{minipage}
\hfill
\begin{minipage}[t]{0.31\linewidth}
\textbf{Contexte géographique}\
\textbf{et socio-économique}
\begin{itemize}\compresslist
\item Densité de population
\item Position socio-économique ("SEP") par commune
%\item Typologie urbain--rural
\item Résultats de votations liées aux politiques familiales
\item Densité des offres de garde d'enfants
\end{itemize}
\end{minipage}

\end{frame}

\begin{frame}{Modèle bayésien: calculer l'excès/déficit de naissance au niveau national en 2000-2025}
\begin{itemize}\compresslist
\item\textbf{Excès/déficit de naissance:} différence entre le nombre de naissances observé et le nombre attendu:\vspace{-0.05cm}
$$E_{a}(t) = y_{a}(t) - \tilde{\mu}_{a}(t)\qquad\qquad E^r_{a}(t) = \dfrac{E_{a}(t)}{\tilde{\mu}_{a}(t)}$$

\vspace{0.1cm}
\item \textbf{Nombre attendu de naissances :} \emph{nombre de naissances si le niveau de fécondité était resté constant au cours des 25 dernières années.}\vspace{0.15cm}
\begin{itemize}\compresslist
\item Évolution du \textbf{nombre de femmes} en âge de procréer et de leur \textbf{structure par âge}
\item \textbf{Report de la maternité} vers des âges plus élevés
\item \textbf{Pas de tendance temporelle de la fécondité}
\end{itemize}

\vspace{0.15cm}
\item\textbf{Modèle bayésien:} Pour chaque âge $a$, on estime le nombre de naissance par mois:
\vspace{0cm}
$$y_{a}(t) \sim \text{Negative-Binomial}\big(\mu_{a}(t),\phi\big)$$
$$\log \mu_{a}(t) =  \textcolor{red}{f}(a+\textcolor{orange}{\theta(t)}) + \textcolor{blue}{\lambda_m(t)} + 
\textcolor{monvert}{\lambda_y(t)} + 
\log N_{a,t}$$
\vspace{0.05cm}
\begin{itemize}\compresslist
\item  \textcolor{red}{Fécondité par âge}: probabilité d'avoir un enfant pour une femme d'âge $a$ (processus Gaussien).
\item  \textcolor{orange}{Report de la maternité} (processus Gaussien)
\item \textbf{\textcolor{blue}{Effet aléatoire mensuel}}
\item \textbf{\textcolor{monvert}{Effet aléatoire annuel}}: corrélation des naissances d'une certaine années
\end{itemize}
\end{itemize}
\end{frame}

\begin{frame}{Estimation de l'excès/déficit de naissances au niveau communal}
\textbf{Désagrégation de la natalité attendue au niveau national}

\begin{itemize}\compresslist
\item L'espérance nationale des naissances est \textbf{répartie entre les communes} selon la structure de la population féminine.
\item Pour chaque mois et groupe d'âge, les naissances attendues sont distribuées entre les $C$ communes selon leur \textbf{part de femmes de cet âge}.
\item Cette désagrégation est réalisée à l'aide d'une \textbf{distribution multinomiale}.
\end{itemize}
$$
\tilde{\mu}_{a,c}(t) \sim \mathrm{Multinomial}\left(\tilde{\mu}_{a}(t),p_{t,a}(1),\ldots, p_{t,a}(C)\right) \qquad \text{avec }p_{t,a}(c) = \dfrac{N_{a,c}(t)}{\sum_c N_{a,c}(t)}
$$
\vspace{0.1cm}

\textbf{Excès de naissances au niveau communal}
$$E_{a,c}(t) = y_{a,c}(t) - \tilde{\mu}_{a,c}(t)\qquad\qquad E^r_{a,c}(t) = \dfrac{E_{a,c}(t)}{\tilde{\mu}_{a,c}(t)}$$
\vspace{0.1cm}

\textbf{Autre niveau de désagrégation}
\begin{itemize}\compresslist
\item \textbf{Nationalité de la mère}: regroupée en 9 régions du monde.%projeter par nationalité
\end{itemize}
\end{frame}


\begin{frame}{Analyses par sous-groupe}

\begin{minipage}[t]{0.55\linewidth}
Le modèle est réestimé en restreignant l'analyse à certains sous-groupes:
\begin{itemize}\compresslist
\item \textbf{Nationalité}
\begin{itemize}\compresslist
\item Mères suisses
\item Mères non suisses
\end{itemize}
\item \textbf{Parité}
\begin{itemize}\compresslist
\item Premières naissances
\item Naissances de rang supérieur
\end{itemize}
\end{itemize}
\vspace*{0.3cm}
\centering
\includegraphics[width=1.1\linewidth,keepaspectratio]{pictures/prop_childless.png}
\end{minipage}
\hfill
\begin{minipage}[t]{0.4\linewidth}
\textbf{Limite pour la parité}

\begin{itemize}\compresslist
\item Dénominateur pas disponible pour la parité: nombre de femmes sans enfant
\item Limite l'interprétation de l'analyse par commune car \emph{proportion de femmes sans enfant} varie géographiquement (p.ex. villes)
\item La \emph{proportion de femmes sans enfant} est estimée empriquement par commune et utilisée pour corriger le dénominateur.
\end{itemize}

\end{minipage}
\end{frame}


\begin{frame}{Calibration du modèle}

\begin{minipage}[h!]{0.34\linewidth}
\textbf{A. Naissances observées vs. attendues}
\begin{itemize}\compresslist
\item Naissances attendues sous \textbf{fécondité constante}
\item Variation observée \textbf{plus grande qu'attendues}
\end{itemize}

\textbf{B. Variation au cours de l'année}
\begin{itemize}\compresslist
\item \textbf{Pic en été}, puis diminution progressive jusqu'à l'hiver
\end{itemize}

\textbf{C--D. Fonction de fécondité par âge}
\begin{itemize}\compresslist
\item Déplacement du \textbf{pic de fécondité} vers la droite.
\item Augmentation de \textbf{2,5 ans en 25 ans}
\end{itemize}
\end{minipage}
\hfill
\begin{minipage}[h!]{0.64\linewidth}
\hspace*{0.1cm}
\centering
\includegraphics[width=1.05\linewidth,keepaspectratio]{pictures/fig1.png}
\end{minipage}
\end{frame}


\begin{frame}{Déficit de naissances à l'échelle nationale}

\begin{minipage}[h!]{0.4\linewidth}

\textbf{A-B. Excès/déficit relatif au cours du temps}

\begin{itemize}\compresslist
\item Déclin quasi continu depuis $\sim$2015
\item En 2025 : \textbf{déficit de 12\%}
\end{itemize}

\textbf{C. Excès de 2021}
\begin{itemize}\compresslist
\item Deux périodes d'excès : \textbf{début d'année et automne}
\item Survenues $\sim$9 mois après les vagues de COVID-19 de 2020
\end{itemize}

\end{minipage}
\hfill
\begin{minipage}[h!]{0.56\linewidth}
\centering
\includegraphics[width=\linewidth,keepaspectratio]{pictures/fig2a.pdf}
\end{minipage}
\end{frame}


\begin{frame}{Déficit des naissances par nationalité }
\begin{minipage}[h!]{0.42\linewidth}
\textbf{A. Déficit moyen entre 2017 et 2025 par région du monde}
\begin{itemize}\compresslist
\item Forte hétérogénéité entre régions
\item Suisse : \textbf{déficit de 10\%}
\item Afrique : \textbf{excès de 62\%}
\end{itemize}
\textbf{B. Évolution au cours du temps}
\begin{itemize}\compresslist
\item Déclin marqué dans \textbf{toutes les régions}
\end{itemize}
\end{minipage}
\hfill
\begin{minipage}[h!]{0.54\linewidth}
\centering
\includegraphics[width=\linewidth,keepaspectratio]{pictures/fig2b.pdf}
\end{minipage}
\end{frame}

\begin{frame}{Variation spatiale du déficit des naissances}
\begin{minipage}[h!]{0.3\linewidth}
\textbf{Excès/déficit relatif moyen entre 2017 et 2025}
\begin{itemize}\compresslist
\item Forte hétérogénéité entre communes
\item Déficit plus élevé dans les \textbf{grandes agglomérations}
\end{itemize}
\textbf{Exceptions régionales}
\begin{itemize}\compresslist
\item Déficit marqué au \textbf{Tessin} et dans certaines parties du \textbf{Valais et des Grisons}
\end{itemize}
\end{minipage}
\hfill
\begin{minipage}[h!]{0.66\linewidth}
\hspace*{-0.6cm}
\centering
\includegraphics[width=1.28\linewidth,keepaspectratio]{pictures/fig3.png}
\end{minipage}
\end{frame}


\begin{frame}{Facteurs contextuels associés au déclin de la natalité}
\begin{minipage}[h!]{0.56\linewidth}
\textbf{Méthode}
\begin{enumerate}\compresslist
\item Les communes sont réparties en \textbf{10 groupes (déciles)} en fonction de la valeur de chaque variable
\item Pour chaque décile, calcul de l'\textbf{excès/déficit moyen} de naissances.
\end{enumerate}
\vspace{0.1cm}
\centering
\includegraphics[width=\linewidth,keepaspectratio]{pictures/fig4.png}
\end{minipage}
\hfill
\begin{minipage}[h!]{0.4\linewidth}
\textbf{Densité de population}
\begin{itemize}\compresslist
\item Excès de naissances dans les communes à \textbf{faible densité}
\item $\sim$\textbf{+5\%} dans le décile le moins dense vs. \textbf{$-$15\%} dans le plus dense
\end{itemize}

\textbf{Position socio-économique}
\begin{itemize}\compresslist
\item Pas d'association monotone claire
\item Excès de naissances aux niveaux de \textbf{SEP très faibles et très élevés}
\end{itemize}
\textbf{Services et politiques familiales}
\begin{itemize}\compresslist
\item Pas d'association positive avec les services de garde ou le soutien à l'extension du congé paternité
\item Associations négatives, probablement liées à la \textbf{densité de population}
\end{itemize}
\end{minipage}
\end{frame}

\begin{frame}{Analyse par sous-groupe: nationalité}

\begin{minipage}[h!]{0.56\linewidth}
\textbf{A. Fécondité}
\begin{itemize}\compresslist
\item Profil de fécondité \textbf{plus précoce et plus prolongé} chez les mères non suisses
\item ICF moyen : \textbf{1,88} vs. \textbf{1,35} enfant par femme
\end{itemize}

\textbf{B. Excès/déficit relatif de naissances}
\begin{itemize}\compresslist
\item Déclin marqué ces dernières années dans les \textbf{deux groupes}
\end{itemize}

\textbf{C. Impact de la densité}
\begin{itemize}\compresslist
\item Association avec la densité \textbf{plus forte chez les mères suisses}
\end{itemize}
\end{minipage}
\hfill
\begin{minipage}[h!]{0.4\linewidth}
\centering
\includegraphics[width=\linewidth,keepaspectratio]{pictures/fig5a.png}
\end{minipage}

\end{frame}

\begin{frame}{Analyse par sous-groupe: parité}

\begin{minipage}[h!]{0.56\linewidth}
\textbf{A. Fécondité}
\begin{itemize}
    \item Fort chevauchement des deux courbes de fécondité
\end{itemize}
\textbf{B. Excès/déficit relatif de naissances}
\begin{itemize}\compresslist
\item Déclin similaire au cours de la dernière décennie
\item Excès de 2021 \textbf{quasi absent pour les premières naissances}
\item Excès de 2021 plus marqué pour les \textbf{naissances d'ordre supérieur}
\end{itemize}

\textbf{C. Impact de la densité}
\begin{itemize}\compresslist
\item Peu de variation selon la densité pour les \textbf{premières naissances}
\item Association négative \textbf{plus forte pour les naissances d'ordre supérieur}
\end{itemize}
\end{minipage}
\hfill
\begin{minipage}[h!]{0.4\linewidth}
\centering
\includegraphics[width=\linewidth,keepaspectratio]{pictures/fig5b.png}
\end{minipage}

\end{frame}



\begin{frame}{Résumé}

\begin{itemize}\compresslist

\item \textbf{Une baisse généralisée des naissances}
  \begin{itemize}\compresslist
  \item Déficit de natalité depuis le milieu des années 2010, malgré le \textbf{report de la maternité}
  \end{itemize}
\vspace{0.2cm}
\item \textbf{Des différences entre les groupes et les régions}
  \begin{itemize}\compresslist
  \item Différences selon la \textbf{nationalité et la parité}, mais une tendance générale à la baisse
  \item Déficits plus marqués dans les \textbf{villes et agglomérations}
  \item Impact de la densité plus marqué chez les \textbf{mères suisses ayant déjà eu un enfant}
  \end{itemize}
\vspace{0.2cm}
\item \textbf{Des mécanismes qui restent à comprendre}
  \begin{itemize}\compresslist
  \item \textbf{Contraintes économiques} : p. ex. niveau socio-économique (SEP)
  \item \textbf{Contraintes pratiques} : p. ex. offre de garde d'enfants
  \item \textbf{Changements de valeurs} : p. ex. préférences concernant la parentalité
  \end{itemize}
  $\Rightarrow$ Données individuelles nécessaires pour étudier ces mécanismes (données hospitalières)
\end{itemize}
\end{frame}

\begin{frame}{Pourquoi la fécondité diminue-t-elle ?}

\begin{minipage}[t]{0.48\linewidth}
\textbf{Biologique}
\begin{itemize}\compresslist
\item Diminution de la \textbf{fertilité}
\item Report de la parentalité vers des âges où la \textbf{fertilité est plus basse}%, dont l'impact est souvent sous-estimé
\end{itemize}
\textbf{Socio-démographique}
\begin{itemize}\compresslist
\item Augmentation de l'\textbf{âge à la parentalité} réduisant la période reproductive et la possibilité d'avoir un \textbf{nombre élevé d'enfants}
\end{itemize}
\vfill
\vspace{0.5cm}
\centering
\includegraphics[width=0.8\linewidth,keepaspectratio]{pictures/fig_desir_enfant.png}

{\tiny Source : DFI.}
\end{minipage}
\hfill
\begin{minipage}[t]{0.48\linewidth}
\textbf{Sociétal / social}
\begin{itemize}\compresslist
\item \textbf{Évolution des valeurs} et individualisation des choix de vie
\item \textbf{Incertitudes face à l'avenir}
\item \textbf{Contraintes financières}
\item Épanouissement \textbf{personnel et professionnel}
\end{itemize}

\vfill
\vspace{0.5cm}
\centering
\includegraphics[width=1\linewidth,keepaspectratio]{pictures/fig_charls2023_raisons.png}
{\tiny Source : CHARLS 2023, enquête nationale suisse sur la fertilité (UZH).}
\end{minipage}
\end{frame}

\begin{frame}{Résumé}

\begin{itemize}\compresslist

\item \textbf{Une baisse généralisée des naissances}
  \begin{itemize}\compresslist
  \item Déficit de natalité depuis le milieu des années 2010, malgré le \textbf{report de la maternité}
  \end{itemize}
\vspace{0.2cm}
\item \textbf{Des différences entre les groupes et les régions}
  \begin{itemize}\compresslist
  \item Différences selon la \textbf{nationalité et la parité}, mais une tendance générale à la baisse
  \item Déficits plus marqués dans les \textbf{villes et agglomérations}
  \item Impact de la densité plus marqué chez les \textbf{mères suisses ayant déjà eu un enfant}
  \end{itemize}
\vspace{0.2cm}
\item \textbf{Des mécanismes qui restent à comprendre}
  \begin{itemize}\compresslist
  \item \textbf{Contraintes économiques} : p. ex. niveau socio-économique (SEP)
  \item \textbf{Contraintes pratiques} : p. ex. offre de garde d'enfants
  \item \textbf{Changements de valeurs} : p. ex. préférences concernant la parentalité
  \end{itemize}
  $\Rightarrow$ Données individuelles nécessaires pour étudier ces mécanismes (données hospitalières)
\end{itemize}
\end{frame}

\end{document}